% GENERATED FILE. Edit canonical lessons, then run npm run generate:books.
% canonical-source-hash: fnv1a64:cecca2c9b68d09fc
% canonical-lessons: ZH-C208-shifu, ZH-C208-jiancha, ZH-C208-anzhuang, ZH-C208-gongju, ZH-C208-dianchi

\chapter{Getting Things Fixed}
\label{ch:zh-getting-things-fixed}
\hlchaptermodality{208}

\begin{chapteropening}
\textbf{By the end of this chapter:} \emph{I can call a repairer, have something checked or put in, and talk about tools and batteries.}

Complete the last lesson of chapter 208: I can call a repairer, have something checked or put in, and talk about tools and batteries.
\end{chapteropening}

\section[shīfu]{shīfu — skilled tradesman; master (and how you address one)}

\label{lesson:ZH-C208-shifu}



\begin{quote}
Before the new one: say the Mandarin for to take a photo, then the Mandarin for tourist sight.
\end{quote}

\subsection*{You'll want to know: shīfu}
\emph{shīfu} — \textquotedblleft{}skilled tradesman; master (and how you address one)\textquotedblright{}.

\begin{quote}
\emph{shīfu wǒ de diànnǎo huài le} — Excuse me, my computer is broken.
\end{quote}

Its characters come later; for now it is a word for the ear.

\begin{cousinweb}
\emph{shīfu} is built from \emph{shī} 'a teacher', the \emph{shī} of \emph{lǎoshī}, and \emph{fu} 'a tutor', both inherited from Middle Chinese, \emph{srij} and \emph{pjuH} (Guangyun, in Baxter's transcription), which glosses the second 'to assist'. In the word the second loses its tone.
\end{cousinweb}

\begin{grammarlens}[title={What you've built}]
One more word for everyday actions.
\end{grammarlens}

\subsection*{Your turn}
\begin{itemize}
\raggedright
  \item \emph{Say it:} \emph{shīfu}
  \item \emph{Say it:} \emph{shīfu}, once more
  \item \emph{Say it:} \emph{shīfu wǒ de diànnǎo huài le} — Excuse me, my computer is broken.
  \item \emph{Recall:} say the Mandarin for to take a photo, then the Mandarin for tourist sight, then say \emph{shīfu} again
\end{itemize}

\begin{tcolorbox}[breakable,colback=teal!4,colframe=teal!35!black,title={Before you move on}]
What does shīfu mean? (\textquotedblleft{}Skilled tradesman; master (and how you address one)\textquotedblright{}.) Say it once more.
\end{tcolorbox}
\section[jiǎnchá]{jiǎnchá — to check; to examine}

\label{lesson:ZH-C208-jiancha}



\begin{quote}
Before the new one: say the Mandarin for tourist sight, then the Mandarin for skilled tradesman.
\end{quote}

\subsection*{You'll want to know: jiǎnchá}
\emph{jiǎnchá} — \textquotedblleft{}to check; to examine\textquotedblright{}.

\begin{quote}
\emph{yīshēng jiǎnchále wǒ de yǎnjing} — The doctor examined my eyes.
\end{quote}

Its characters come later; for now it is a word for the ear.

\begin{cousinweb}
\emph{jiǎnchá} is built from two words for looking into things, both inherited from Middle Chinese, \emph{kjemX} and \emph{dzrae} (Guangyun, in Baxter's transcription). The dictionary glosses the first as the seal on a document and 'to check'; it glosses the second as a raft, so how it came to mean examining is not settled.
\end{cousinweb}

\begin{grammarlens}[title={What you've built}]
One more word for everyday actions.
\end{grammarlens}

\subsection*{Your turn}
\begin{itemize}
\raggedright
  \item \emph{Say it:} \emph{jiǎnchá}
  \item \emph{Say it:} \emph{jiǎnchá}, once more
  \item \emph{Say it:} \emph{yīshēng jiǎnchále wǒ de yǎnjing} — The doctor examined my eyes.
  \item \emph{Recall:} say the Mandarin for tourist sight, then the Mandarin for skilled tradesman, then say \emph{jiǎnchá} again
\end{itemize}

\begin{tcolorbox}[breakable,colback=teal!4,colframe=teal!35!black,title={Before you move on}]
What does jiǎnchá mean? (\textquotedblleft{}To check; to examine\textquotedblright{}.) Say it once more.
\end{tcolorbox}
\section[ānzhuāng]{ānzhuāng — to install; to put in}

\label{lesson:ZH-C208-anzhuang}



\begin{quote}
Before the new one: say the Mandarin for skilled tradesman, then the Mandarin for to examine.
\end{quote}

\subsection*{You'll want to know: ānzhuāng}
\emph{ānzhuāng} — \textquotedblleft{}to install; to put in\textquotedblright{}.

\begin{quote}
\emph{tāmen zài ānzhuāng diànshì} — They're installing the television.
\end{quote}

Its characters come later; for now it is a word for the ear.

\begin{cousinweb}
\emph{ānzhuāng} is built from \emph{ān} 'to settle' and \emph{zhuāng} 'to fit out', both inherited from Middle Chinese, \emph{'an} and \emph{tsrjang} (Guangyun, in Baxter's transcription), which glosses the first 'calm' and the second 'to pack up' and 'to dress'. Putting a machine in place is not given there, so how the pair came to it is not settled.
\end{cousinweb}

\begin{grammarlens}[title={What you've built}]
One more word for everyday actions.
\end{grammarlens}

\subsection*{Your turn}
\begin{itemize}
\raggedright
  \item \emph{Say it:} \emph{ānzhuāng}
  \item \emph{Say it:} \emph{ānzhuāng}, once more
  \item \emph{Say it:} \emph{tāmen zài ānzhuāng diànshì} — They're installing the television.
  \item \emph{Recall:} say the Mandarin for skilled tradesman, then the Mandarin for to examine, then say \emph{ānzhuāng} again
\end{itemize}

\begin{tcolorbox}[breakable,colback=teal!4,colframe=teal!35!black,title={Before you move on}]
What does ānzhuāng mean? (\textquotedblleft{}To install; to put in\textquotedblright{}.) Say it once more.
\end{tcolorbox}
\section[gōngjù]{gōngjù — tool}

\label{lesson:ZH-C208-gongju}



\begin{quote}
Before the new one: say the Mandarin for to examine, then the Mandarin for to install.
\end{quote}

\subsection*{You'll want to know: gōngjù}
\emph{gōngjù} — \textquotedblleft{}tool\textquotedblright{}.

\begin{quote}
\emph{shīfu dàile gōngjù} — The repairman brought his tools.
\end{quote}

Its characters come later; for now it is a word for the ear.

\begin{cousinweb}
\emph{gōngjù} is built from \emph{gōng} 'work' and \emph{jù} 'an implement', both inherited from Middle Chinese, \emph{kuwng} and \emph{gjuH} (Guangyun, in Baxter's transcription), which glosses the second 'to prepare'.
\end{cousinweb}

\begin{grammarlens}[title={What you've built}]
One more word for everyday actions.
\end{grammarlens}

\subsection*{Your turn}
\begin{itemize}
\raggedright
  \item \emph{Say it:} \emph{gōngjù}
  \item \emph{Say it:} \emph{gōngjù}, once more
  \item \emph{Say it:} \emph{shīfu dàile gōngjù} — The repairman brought his tools.
  \item \emph{Recall:} say the Mandarin for to examine, then the Mandarin for to install, then say \emph{gōngjù} again
\end{itemize}

\begin{tcolorbox}[breakable,colback=teal!4,colframe=teal!35!black,title={Before you move on}]
What does gōngjù mean? (\textquotedblleft{}Tool\textquotedblright{}.) Say it once more.
\end{tcolorbox}
\section[diànchí]{diànchí — battery}

\label{lesson:ZH-C208-dianchi}



\begin{quote}
Before the new one: say the Mandarin for to install, then the Mandarin for tool.
\end{quote}

\subsection*{You'll want to know: diànchí}
\emph{diànchí} — \textquotedblleft{}battery\textquotedblright{}.

\begin{quote}
\emph{shǒujī de diànchí huài le} — The phone's battery has gone.
\end{quote}

Its characters come later; for now it is a word for the ear.

\begin{cousinweb}
\emph{diànchí} is built from \zh{电} 'electricity' and \emph{chí} 'a pool', both inherited from Middle Chinese, \emph{denH} and \emph{drje} (Guangyun, in Baxter's transcription), which glosses the second as standing water: a pool of electricity.
\end{cousinweb}

\begin{grammarlens}[title={What you've built}]
One more word for everyday actions.
\end{grammarlens}

\subsection*{Your turn}
\begin{itemize}
\raggedright
  \item \emph{Say it:} \emph{diànchí}
  \item \emph{Say it:} \emph{diànchí}, once more
  \item \emph{Say it:} \emph{shǒujī de diànchí huài le} — The phone's battery has gone.
  \item \emph{Recall:} say the Mandarin for to install, then the Mandarin for tool, then say \emph{diànchí} again
\end{itemize}

\begin{tcolorbox}[breakable,colback=teal!4,colframe=teal!35!black,title={Before you move on}]
What does diànchí mean? (\textquotedblleft{}Battery\textquotedblright{}.) Say it once more.
\end{tcolorbox}
